% !PRINTED 404
\weberpage{404}
Wir können aber auch noch weiter annehmen, dass $a_0 = 1$ sei; denn setzen wir in dem Product $a_0^{n-1} f(x)$
\[
a_0 x = x_1,
\]
so wird die Gleichung
\[
x_1^n + a_1 x_1^{n-1} + a_0 a_2 x_1^{n-2} + \cdots + a_0^{n-1} a_n = 0.
\]
Wir wollen also jetzt annehmen, dass in der Gleichung
\weq{1}{x^n + a_1 x^{n-1} + a_2 x^{n-2} + \cdots + a_{n-1} x + a_n = 0}
die Co\"efficienten $a_1, a_2 \ldots a_n$ ganze Zahlen sind. Eine \gesperrt{rationale} Wurzel von (1) kann nicht eine gebrochene Zahl sein; denn nehmen wir an, es werde (1) befriedigt durch den Bruch
\[
x = \frac{p}{q},
\]
wo $p$ und $q$ ganze Zahlen ohne gemeinsamen Theiler sind, und $q$ positiv und grösser als 1 ist, so folgt durch Multiplication mit $q^{n-1}$
\[
\frac{p^n}{q} + a_1 p^{n-1} + a_2 q p^{n-2} + \cdots + a_n q^{n-1} = 0.
\]
Es müsste also $p^n : q$ eine ganze Zahl sein, was unmöglich ist.

Wenn nun die Gleichung (1) dadurch befriedigt werden kann, dass man für $x$ eine ganze Zahl setzt, so muss diese ganze Zahl, wie die Gleichung (1) zeigt, nothwendig ein Theiler von $a_n$ sein, und man hat also nur die verschiedenen Divisoren von $a_n$, mit positivem und negativem Zeichen behaftet, versuchsweise in die Gleichung einzusetzen.

Auf die Frage nach rationalen Wurzeln einer ganzzahligen Gleichung kommt auch die allgemeine Frage zurück, ob eine ganze rationale Function
\[
f(x) = x^n + a_1 x^{n-1} + a_2 x^{n-2} + \cdots + a_{n-1} x + a_n,
\]
deren Co\"efficienten $a_1, a_2 \ldots a_n$ ganze Zahlen sind, durch eine rationale Function von niedrigerem Grade $\nu$:
\[
\varphi(x) = x^\nu + \alpha_1 x^{\nu-1} + \alpha_2 x^{\nu-2} + \cdots + \alpha_{\nu-1} x + \alpha_\nu
\]
mit rationalen Co\"fficienten $\alpha_1, \alpha_2 \ldots \alpha_\nu$ theilbar sein kann. Wenn $f(x)$ durch $\varphi(x)$ theilbar ist, so ist der Quotient $\psi(x)$ wieder eine ganze rationale Function mit rationalen Co\"fficienten vom Grade $n - \nu$:
\[
\psi(x) = x^{n-\nu} + \beta_1 x^{n-\nu-1} + \cdots + \beta_{n-\nu-1} x + \beta_{n-\nu},
\]
und da $f(x)$ auch durch $\psi(x)$ theilbar ist, so können wir bei
%CONT>>
